% ##############################################################################
\section{NoesisServer: fulfillment and monitoring}
\label{sec:noesis_server}

\NoesisServer{} is the layer that turns a matched contract
(Definition~\ref{def:contract}) into actual inference requests against real
providers, and turns the resulting traffic into a measurement of whether the
contract was honored and the quality of its fulfillment.

% ==============================================================================
\subsection{Gateway architecture}
\label{sec:gateway_architecture}

\NoesisServer{} is a lightweight LLM API gateway modeled on
multi-provider routers (e.g., OpenRouter): it exposes a single API,
proxies each request to one or more underlying providers, and streams the
response back to the caller. Two properties distinguish it from a plain
pass-through router. First, every request is tagged with the contract
$\kappa$ it is served under, if any, so that fulfillment can be attributed
(Section~\ref{sec:fulfillment_monitoring}). Second, every prompt/response
pair is logged with its provider, model, measured latency and throughput,
measured cost,
and, when available, a user feedback signal, regardless of whether the
request was placed under a contract. The storage schema is intentionally
minimal:

$$
  (\text{prompt}, \text{response}, \text{provider}, \text{model},
   \text{latency}, \text{throughput}, \text{cost}, \text{contract})
$$

This logging is passive by default: it does not change routing behavior and
exists purely to build a reusable dataset for evaluation. Privacy and
data-handling safeguards (PII scrubbing, opt-in/opt-out) are necessary before
logging real traffic and are discussed as an open question in
Section~\ref{sec:open_questions}.

Figure~\ref{fig:request_path} traces this path end to end: the caller's
request is proxied through a router that selects one of several
candidate providers, the response is streamed back, and on the way back,
the request is logged as the tuple (prompt, response, provider, model, latency,
throughput, cost, contract) described above.

% rendered_images:begin
% ```graphviz
% digraph RequestPath {
%   // ---------------------------------------------------------------
%   bgcolor="transparent";
%   pad="0.10";
%   splines=spline;
%   nodesep=0.32;
%   ranksep=0.60;
% 
%   node [shape=box,
%         style="rounded,filled",
%         fillcolor="#FFFFFF",
%         color="#9AA9B8",
%         penwidth=1.8,
%         fontname="Helvetica",
%         fontcolor="#243B53",
%         fontsize=10,
%         margin="0.20,0.12",
%         height=0.45];
% 
%   edge [color="#A3B1C0",
%         penwidth=1.3,
%         arrowhead=vee,
%         arrowsize=0.75,
%         fontname="Helvetica",
%         fontsize=8.5,
%         fontcolor="#7B8794"];
% 
%   // ---------------------------------------------------------------
%   rankdir=LR;
% 
%   // actor : rose   |   process : blue   |   external : violet   |   artifact : sage green
% 
%   Caller  [label=<<b>Caller</b><br/><font point-size="8.5" color="#A34F6C">Application / agent</font>>,
%            fillcolor="#F6E1E8", color="#D98CA8", fontcolor="#6B2A44"];
%   Gateway [label=<<b>API gateway</b><br/><font point-size="8.5" color="#41719C">NoesisServer</font>>,
%            fillcolor="#D3E3F3", color="#7CA6CE", fontcolor="#1F4E79"];
%   Router  [shape=invtrapezium, label=<<b>Route by κ</b><br/><font point-size="7.5" color="#41719C">1-of-3 select</font>>,
%            fillcolor="#D3E3F3", color="#7CA6CE", fontcolor="#1F4E79", height=0.4, margin="0.10,0.02"];
% 
%   ProviderA [label="Provider A", fillcolor="#E9E1F7", color="#A88FD9", fontcolor="#45296B"];
%   ProviderB [label="Provider B", fillcolor="#E9E1F7", color="#A88FD9", fontcolor="#45296B"];
%   ProviderC [label="Provider C", fillcolor="#E9E1F7", color="#A88FD9", fontcolor="#45296B"];
% 
%   Log [label=<<font point-size="8.5">Log:<br/>(prompt, response, provider, model,<br/>latency, throughput, cost, contract-or-bottom)</font>>,
%        shape=note, fillcolor="#DFEDE0", color="#8FB79A", fontcolor="#2E5A3D"];
% 
%   Caller -> Gateway [label="request", penwidth=1.6, color="#8592A3"];
%   Gateway -> Router [label="request", penwidth=1.6, color="#8592A3"];
%   Router -> ProviderA;
%   Router -> ProviderB;
%   Router -> ProviderC;
%   Router -> Gateway [label="response", style=dashed, constraint=false];
%   Gateway -> Caller [label="response", style=dashed];
%   Gateway -> Log [
%     label="logged on the way back",
%     style=dashed, penwidth=1.5,
%     color="#C0455B", fontcolor="#C0455B",
%     constraint=false
%   ];
% }
% ```
% label=fig:request_path
% caption=Path of a single request through \NoesisServer{}, routed by a demux-style selector to one of several candidate providers.
% rendered_images:end
% render_images:begin
\begin{figure}[H]
  \includegraphics[width=\linewidth]{05_noesis_server.tex.figs/05_noesis_server.1.png}
  \caption{Path of a single request through \NoesisServer{}, routed by a demux-style selector to one of several candidate providers.}
  \label{fig:request_path}
\end{figure}
% render_images:end

% ==============================================================================
\subsection{Contract fulfillment monitoring}
\label{sec:fulfillment_monitoring}

For a request $x$ served under contract $\kappa$, let
$(c, \ell, \eta) = \sclassattr(\kappa)$ be the contracted service class. Let
$\hat c(x) \in \calC$ denote the capability tier actually delivered,
estimated from the response (e.g., via a verifier model, a benchmark probe,
or agreement with a stronger reference model), and let $\hat L(x)$ and
$\hat H(x)$ denote the measured latency and output tokens per second. The
capability estimator is the pluggable capability-measurement component of
Section~\ref{sec:modularity}: the fulfillment test below only requires
$\hat c(x)$, not a commitment to one specific way of computing it. All three
dimensions are checked per request, and request $x$ is a \emph{success} under
$\kappa$ only if it meets all of them:

\begin{equation}
  \label{eq:success_indicator}
  s(x) = \mathbb{1}\!\left[
    \hat c(x) \succeq c
    \;\wedge\;
    \hat L(x) \leq L_{\max}(\ell)
    \;\wedge\;
    \hat H(x) \geq H_{\min}(\eta)
  \right].
\end{equation}

A request that is not a success is one \emph{violation}, however many
dimensions it misses.

\begin{definition}[Measured reliability]
\label{def:measured_reliability}
Let $W$ denote the set of requests served under contract $\kappa$ over the
life of the contract. The measured reliability, one minus the fraction of
requests that are violations, is
\begin{equation}
  \label{eq:measured_reliability}
  \hat r(\kappa, W) = \frac{1}{|W|} \sum_{x \in W} s(x).
\end{equation}
\end{definition}

A naive violation rule, flagging $\kappa$ whenever
$\hat r(\kappa, W) < \rminattr(\kappa)$, conflates genuine under-delivery
with sampling noise: $\hat r(\kappa, W)$ is a sample average of an
unobserved true success rate, and a small $|W|$ makes false accusations
likely. We instead treat fulfillment monitoring as a one-sided hypothesis
test: the null hypothesis $H_0$ is that the seller is compliant, i.e., its
true success rate is at least $\rminattr(\kappa)$, and $\kappa$ is flagged
as violated only if $H_0$ is rejected at a chosen confidence level
$1 - \delta$. A simple normal-approximation upper confidence bound gives

\begin{equation}
  \label{eq:reliability_upper_bound}
  \hat r_{\text{upper}}(\kappa, W)
  = \hat r(\kappa, W)
  + z_{1 - \delta}
    \sqrt{
      \frac{\hat r(\kappa, W) \, (1 - \hat r(\kappa, W))}{|W|}
    },
\end{equation}
where $z_{1 - \delta}$ is the corresponding standard normal quantile, and
$\kappa$ is flagged as violated when
$\hat r_{\text{upper}}(\kappa, W) < \rminattr(\kappa)$. For small $|W|$, an
exact interval (e.g., Clopper--Pearson) is preferable to the normal
approximation used here for exposition. Flagged violations are reported
to \NoesisMarket{} and feed the reputation update of
Equation~\eqref{eq:reputation_update}.

\paragraph{Capability measured on a sample.}
Latency and throughput are observed on every request, but capability may be
observable only on a sample, e.g., known-answer probes that the platform sends
under the contract. Pooling the probes with all other requests would dilute a
capability shortfall, since most requests then carry no capability evidence.
Instead, since the probability of a success is the probability of being served
within the limits times the probability of being capable given that,
reliability is estimated as a product of two factors,
\begin{equation}
  \label{eq:product_reliability}
  \hat r(\kappa, W) = \hat r_{\text{serve}} \cdot \hat r_{\text{cap}},
  \qquad
  \hat r_{\text{upper}}(\kappa, W) = U_{\text{serve}} \cdot U_{\text{cap}}.
\end{equation}
Here $\hat r_{\text{serve}}$ is the fraction of all requests in $W$, probes
included, answered within $L_{\max}(\ell)$ and at $H_{\min}(\eta)$ or more;
$\hat r_{\text{cap}}$ is the fraction of the \emph{answered} probes that reach
capability $c$; and $U_{\text{serve}}$ and $U_{\text{cap}}$ are the upper
confidence bounds of the two factors at level $1 - \delta$ (e.g., exact
Clopper--Pearson bounds). A timeout or an error therefore counts once, in
$\hat r_{\text{serve}}$, and not also as a wrong answer. The violation test,
the reputation update of Equation~\eqref{eq:reputation_update}, and slashing
(Definition~\ref{def:staked_ask}) use this $\hat r_{\text{upper}}$ unchanged.
Since the product falls below the true success probability only if one of the
two bounds falls below its factor, a compliant seller is flagged with
probability at most about $2\delta$. The estimate assumes that answered probes
are as likely to reach capability $c$ as buyer requests served within the
limits, which fails if a seller can recognize probes.

\begin{example}[Fulfillment reporting]
\label{ex:fulfillment_reporting}
A contract promising 99.9\% reliability at class
(\texttt{high}, \texttt{fast}, \texttt{medium}) is
served by three providers over a window of 5,000 requests. \NoesisServer{}
computes $\hat r(\kappa, W) = 0.994$ and, applying
Equation~\eqref{eq:reliability_upper_bound} with $\delta = 0.05$, finds
$\hat r_{\text{upper}}(\kappa, W) \approx 0.996 < 0.999$: the shortfall is
unlikely to be sampling noise, so the contract is flagged as violated and reported back
to \NoesisMarket.
\end{example}

Figure~\ref{fig:reliability_test} places the point estimate, its upper
confidence bound, and the threshold on a single number line for both
outcomes of the test: when the interval
$[\hat r(\kappa, W), \hat r_{\text{upper}}]$ reaches $\rminattr(\kappa)$, the
shortfall, if any, is attributed to sampling noise, and when the whole
interval lies below $\rminattr(\kappa)$, $\kappa$ is flagged as violated and
reported to \NoesisMarket.

% rendered_images:begin
% ```tikz
% \definecolor{flagzone}{RGB}{248,220,220}
% \definecolor{flagline}{RGB}{107,31,31}
% \definecolor{okzone}{RGB}{223,237,224}
% \definecolor{okline}{RGB}{46,90,61}
% \definecolor{dividercolor}{RGB}{133,146,163}
% \definecolor{markcolor}{RGB}{36,59,83}
% \fontfamily{phv}\selectfont
% \fill[flagzone, opacity=0.45] (0,-1) rectangle (6,3.4);
% \fill[okzone, opacity=0.45] (6,-1) rectangle (12,3.4);
% \draw[dashed, thick, dividercolor] (6,-1) -- (6,3.4);
% \node[above] at (6,3.4) {\footnotesize $R_{\min}(\kappa)$};
% \draw[very thick, okline] (4.6,2.4) -- (8.4,2.4);
% \draw[very thick, okline] (4.6,2.25) -- (4.6,2.55);
% \draw[very thick, okline] (8.4,2.25) -- (8.4,2.55);
% \filldraw[markcolor] (4.6,2.4) circle (2pt);
% \node[below] at (8.4,2.1) {\footnotesize $\hat r_{\text{upper}}$};
% \node[above] at (4.6,2.7) {\footnotesize $\hat r(\kappa,W)$};
% \node[align=right, okline] at (-1.9,2.4) {\footnotesize compliant:\\[-1pt]\footnotesize attributed to noise};
% \draw[very thick, flagline] (1.6,0) -- (4.8,0);
% \draw[very thick, flagline] (1.6,-0.15) -- (1.6,0.15);
% \draw[very thick, flagline] (4.8,-0.15) -- (4.8,0.15);
% \filldraw[markcolor] (1.6,0) circle (2pt);
% \node[below] at (4.8,-0.3) {\footnotesize $\hat r_{\text{upper}}$};
% \node[above] at (1.6,0.3) {\footnotesize $\hat r(\kappa,W)$};
% \node[align=right, flagline] at (-1.9,0) {\footnotesize flagged:\\[-1pt]\footnotesize violation reported};
% ```
% label=fig:reliability_test
% caption=The reliability hypothesis test on a single number line.
% rendered_images:end
% render_images:begin
\begin{figure}[H]
  \includegraphics[width=\linewidth]{05_noesis_server.tex.figs/05_noesis_server.2.png}
  \caption{The reliability hypothesis test on a single number line.}
  \label{fig:reliability_test}
\end{figure}
% render_images:end

% ==============================================================================
\subsection{Difficulty-aware routing}
\label{sec:routing}

Not every request needs the strongest available model. \NoesisServer{}
trains a difficulty estimator
\begin{equation}
  \label{eq:difficulty_estimator}
  \hat d : \calP \to \calC
\end{equation}
where $\calP$ denotes the space of prompts, mapping a prompt $x \in \calP$
to the cheapest capability tier expected to
answer it correctly, and routes each request to the tier the estimator
itself outputs,
\begin{equation}
  \label{eq:routing_policy}
  \pi(x) = \hat d(x).
\end{equation}

Routing is only worthwhile if the estimator is cheaper than the savings it
produces. Let $\text{cost}(c)$ denote the per-request cost of serving tier
$c$, non-decreasing in $c$, let $c_K$ denote the strongest tier, and let
$c_{\hat d}$ denote the per-request cost of evaluating $\hat d$ itself.
Routing yields a net saving over an always-route-to-$c_K$ baseline exactly
when

\begin{equation}
  \label{eq:routing_saving_condition}
  c_{\hat d} + \EE_x [\, \text{cost}(\pi(x)) \,]
  <
  \text{cost}(c_K).
\end{equation}

Inequality~\eqref{eq:routing_saving_condition} makes explicit the trap
flagged in the underlying design notes: an estimator that is itself
expensive, or inaccurate enough to route hard requests to weak tiers and
force costly retries, can cost more than always calling the strong model.
Measuring the left-hand side against the right-hand side on real traffic
is a prerequisite for deploying routing in production, and is deferred to
future implementation work (Section~\ref{sec:open_questions}).

Figure~\ref{fig:routing_flowchart} lays out this chain next to the
condition under which it pays off: a prompt $x$ is mapped by the
difficulty estimator to the cheapest capability tier $\pi(x)$ expected to
answer it correctly, and routing only pays off when the estimator's own
cost plus the expected routed cost undercuts always calling the strongest
tier (Inequality~\eqref{eq:routing_saving_condition}).

% rendered_images:begin
% ```graphviz
% digraph Routing {
%   // ---------------------------------------------------------------
%   bgcolor="transparent";
%   pad="0.06";
%   splines=spline;
%   nodesep=0.45;
%   ranksep=0.60;
% 
%   node [shape=box,
%         style="rounded,filled",
%         fillcolor="#FFFFFF",
%         color="#9AA9B8",
%         penwidth=1.8,
%         fontname="Helvetica",
%         fontcolor="#243B53",
%         fontsize=10,
%         margin="0.22,0.14",
%         height=0.50];
% 
%   edge [color="#A3B1C0",
%         penwidth=1.3,
%         arrowhead=vee,
%         arrowsize=0.75,
%         fontname="Helvetica",
%         fontsize=8.5,
%         fontcolor="#7B8794"];
% 
%   // ---------------------------------------------------------------
%   rankdir=LR;
% 
%   // source : warm sand   |   process : blue   |   note : violet
% 
%   Prompt    [label=<<b>Prompt x</b>>,
%              fillcolor="#FBEBD4", color="#D9A85F", fontcolor="#6B4517"];
%   Estimator [label=<<b>Difficulty estimator</b><br/><font point-size="8.5" color="#41719C">Cheapest tier expected<br/>to answer correctly</font>>,
%              fillcolor="#D3E3F3", color="#7CA6CE", fontcolor="#1F4E79"];
%   Tier      [label=<<b>Routed tier π(x)</b><br/><font point-size="8.5" color="#41719C">Cheapest sufficient tier</font>>,
%              fillcolor="#D3E3F3", color="#7CA6CE", fontcolor="#1F4E79"];
%   Cost      [label=<<font point-size="8.5">Net saving only if:<br/>estimator cost + expected routed cost<br/>&lt; always-strong-model cost</font>>,
%              shape=note, fillcolor="#E9E1F7", color="#A88FD9", fontcolor="#45296B"];
% 
%   Prompt -> Estimator [penwidth=1.6, color="#8592A3"];
%   Estimator -> Tier [penwidth=1.6, color="#8592A3"];
%   Tier -> Cost [style=dashed, constraint=false];
% }
% ```
% label=fig:routing_flowchart
% caption=Difficulty-aware routing to the cheapest sufficient capability tier.
% rendered_images:end
% render_images:begin
\begin{figure}[H]
  \includegraphics[width=\linewidth]{05_noesis_server.tex.figs/05_noesis_server.3.png}
  \caption{Difficulty-aware routing to the cheapest sufficient capability tier.}
  \label{fig:routing_flowchart}
\end{figure}
% render_images:end

% ==============================================================================
\subsection{Answer fusion}
\label{sec:fusion}

Fusion is the complementary lever to routing. Instead of picking one tier
per request, \NoesisServer{} may query a set of providers $S(x) \subseteq
\calM$, where $\calM$ denotes the set of available providers, and combine
their answers through an aggregator
\begin{equation}
  \label{eq:fusion_aggregator}
  \hat y(x) = g\big( \{ f_m(x) \}_{m \in S(x)} \big),
\end{equation}
where $g$ may be a majority vote, a verifier model that scores and selects
among candidates, or a learned combiner trained on logged outcomes: the
aggregator is the pluggable answer-fusion component of
Section~\ref{sec:modularity}.

Routing and fusion trade cost for accuracy in opposite directions relative
to a single-strong-model baseline: routing accepts a higher risk of a
lower-quality answer on some fraction of requests in exchange for a lower
expected cost, while fusion accepts a higher expected cost, from querying
multiple providers per request, in exchange for a lower risk of a
low-quality answer. Which lever is preferable for a given contract depends
on the contract's price versus its $\rminattr$ term: a low-price,
low-$\rminattr$ contract favors routing, while a high-$\rminattr$ contract
may only be economically served through fusion.

Figure~\ref{fig:fusion_fanout} contrasts fusion's fan-out/fan-in shape
with routing's single-provider path: fusion sends the same prompt to
every provider in $S(x)$ and combines their answers through the
aggregator $g$ (Equation~\eqref{eq:fusion_aggregator}), trading extra
inference cost for a lower risk of a low-quality answer.

% rendered_images:begin
% ```graphviz
% digraph Fusion {
%   // ---------------------------------------------------------------
%   bgcolor="transparent";
%   pad="0.06";
%   splines=spline;
%   nodesep=0.30;
%   ranksep=0.46;
% 
%   node [shape=box,
%         style="rounded,filled",
%         fillcolor="#FFFFFF",
%         color="#9AA9B8",
%         penwidth=1.8,
%         fontname="Helvetica",
%         fontcolor="#243B53",
%         fontsize=9.5,
%         margin="0.18,0.10",
%         height=0.45];
% 
%   edge [color="#A3B1C0",
%         penwidth=1.3,
%         arrowhead=vee,
%         arrowsize=0.75,
%         fontname="Helvetica",
%         fontsize=8.5,
%         fontcolor="#7B8794"];
% 
%   // ---------------------------------------------------------------
%   rankdir=LR;
% 
%   // source : warm sand   |   external : violet   |   process : blue   |   artifact : sage green
% 
%   subgraph cluster_routing {
%     label     = "Routing (single path)";
%     labelloc  = "t";
%     fontname  = "Helvetica-Bold";
%     fontsize  = 10.5;
%     fontcolor = "#3A4A5C";
%     style     = "rounded,filled";
%     fillcolor = "#F7F9FB";
%     color     = "#C7D0DA";
%     penwidth  = 1.0;
%     margin    = 14;
% 
%     PromptR   [label=<<b>prompt x</b>>,
%                fillcolor="#FBEBD4", color="#D9A85F", fontcolor="#6B4517"];
%     ProviderR [label=<<b>one provider</b><br/><font point-size="8" color="#6B52A8">routed tier</font>>,
%                fillcolor="#E9E1F7", color="#A88FD9", fontcolor="#45296B"];
%     AnswerR   [label=<<b>answer</b>>,
%                fillcolor="#DFEDE0", color="#8FB79A", fontcolor="#2E5A3D"];
%     PromptR -> ProviderR -> AnswerR;
%   }
% 
%   subgraph cluster_fusion {
%     label     = "Fusion (fan-out / fan-in)";
%     labelloc  = "t";
%     fontname  = "Helvetica-Bold";
%     fontsize  = 10.5;
%     fontcolor = "#3A4A5C";
%     style     = "rounded,filled";
%     fillcolor = "#F7F9FB";
%     color     = "#C7D0DA";
%     penwidth  = 1.0;
%     margin    = 14;
% 
%     PromptF    [label=<<b>prompt x</b>>,
%                 fillcolor="#FBEBD4", color="#D9A85F", fontcolor="#6B4517"];
%     P1         [label=<<b>provider in S(x)</b>>,
%                 fillcolor="#E9E1F7", color="#A88FD9", fontcolor="#45296B"];
%     P2         [label=<<b>provider in S(x)</b>>,
%                 fillcolor="#E9E1F7", color="#A88FD9", fontcolor="#45296B"];
%     P3         [label=<<b>provider in S(x)</b>>,
%                 fillcolor="#E9E1F7", color="#A88FD9", fontcolor="#45296B"];
%     Aggregator [label=<<b>aggregator g</b><br/><font point-size="8" color="#41719C">vote / verifier / learned</font>>,
%                 fillcolor="#D3E3F3", color="#7CA6CE", fontcolor="#1F4E79"];
%     AnswerF    [label=<<b>answer</b>>,
%                 fillcolor="#DFEDE0", color="#8FB79A", fontcolor="#2E5A3D"];
%     PromptF -> P1;
%     PromptF -> P2;
%     PromptF -> P3;
%     P1 -> Aggregator;
%     P2 -> Aggregator;
%     P3 -> Aggregator;
%     Aggregator -> AnswerF;
%   }
% }
% ```
% label=fig:fusion_fanout
% caption=Answer fusion's fan-out/fan-in versus routing's single-provider path.
% rendered_images:end
% render_images:begin
\begin{figure}[H]
  \includegraphics[width=\linewidth]{05_noesis_server.tex.figs/05_noesis_server.4.png}
  \caption{Answer fusion's fan-out/fan-in versus routing's single-provider path.}
  \label{fig:fusion_fanout}
\end{figure}
% render_images:end

% ==============================================================================
\subsection{Training signal and opt-in distillation}
\label{sec:distillation}

The difficulty estimator of Equation~\eqref{eq:difficulty_estimator}
needs a training signal, which the logged traffic supplies for free.

\begin{definition}[Agreement label]
\label{def:agreement_label}
For a prompt $x$ answered by both a cheap model $f_{\text{cheap}}$ and a
strong model $f_{\text{strong}}$, the agreement label is
\begin{equation}
  \label{eq:agreement_label}
  y(x) = 1 - \mathbb{1}\big[ f_{\text{cheap}}(x) \equiv f_{\text{strong}}(x) \big],
\end{equation}
where $\equiv$ denotes exact match or a similarity threshold, used as a
proxy label for ``this prompt was hard'' ($y(x) = 1$) versus ``easy''
($y(x) = 0$).
\end{definition}

Definition~\ref{def:agreement_label} is a proxy, not a ground-truth
correctness label: the two models can agree on a wrong answer, or
disagree while the cheap model happens to be correct. A genuine quality
label requires independent verification, for instance from a pairwise
preference arena, which is a companion research direction discussed in
Section~\ref{sec:related_work}. With that caveat, the logged corpus
$$
  \calD = \{ (x_i, f_{\text{strong}}(x_i)) \}_{i=1}^N
$$
also supports knowledge distillation of a smaller model $f_\theta$ trained
to approximate $f_{\text{strong}}$'s outputs on the observed task
distribution \cite{hinton2015}, offering a second, independent way for
\NoesisServer{} to reduce cost on traffic it has already observed.

Unlike the passive logging of Section~\ref{sec:gateway_architecture},
inclusion in the distillation corpus $\calD$ is an \emph{opt-in
distillation program}: a customer's traffic contributes to $\calD$ only
after the customer opts in, and only after prompts and responses are
anonymized to remove customer- and end-user-identifying content. Because
an opted-in customer's traffic makes the corpus, and therefore the
distilled model, more useful for every buyer on \NoesisMarket, an
opted-in customer can be offered a preferential rate in exchange, in the
same spirit as data-sharing discounts in other metered-service markets.

Figure~\ref{fig:training_pipeline} shows the two consumers that share
this one logged dataset: logged traffic feeds agreement labels
(Definition~\ref{def:agreement_label}), which train the difficulty
estimator of Section~\ref{sec:routing}, while the same traffic also forms
the corpus $\calD$ used to distill a smaller model $f_\theta$
(Section~\ref{sec:distillation}), as two parallel, independent consumers
of one dataset.

% rendered_images:begin
% ```graphviz
% digraph TrainingPipeline {
%   // ---------------------------------------------------------------
%   bgcolor="transparent";
%   pad="0.06";
%   splines=spline;
%   nodesep=0.30;
%   ranksep=0.46;
% 
%   node [shape=box,
%         style="rounded,filled",
%         fillcolor="#FFFFFF",
%         color="#9AA9B8",
%         penwidth=1.8,
%         fontname="Helvetica",
%         fontcolor="#243B53",
%         fontsize=10,
%         margin="0.22,0.14",
%         height=0.50];
% 
%   edge [color="#A3B1C0",
%         penwidth=1.3,
%         arrowhead=vee,
%         arrowsize=0.75,
%         fontname="Helvetica",
%         fontsize=8.5,
%         fontcolor="#7B8794"];
% 
%   // ---------------------------------------------------------------
%   rankdir=LR;
% 
%   // source  : warm sand   |   artifact : sage green   |   process : blue
% 
%   Traffic   [label=<<b>Logged traffic</b><br/><font point-size="8.5" color="#946A2E">Prompt / response pairs</font>>,
%              fillcolor="#FBEBD4", color="#D9A85F", fontcolor="#6B4517"];
% 
%   Agreement [label=<<b>Agreement labels</b><br/><font point-size="8.5" color="#4F7A5A">Cheap vs. strong model</font>>,
%              fillcolor="#DFEDE0", color="#8FB79A", fontcolor="#2E5A3D"];
% 
%   Corpus    [label=<<b>Corpus <i>D</i></b><br/><font point-size="8.5" color="#4F7A5A">(x, strong&#8209;model answer)</font>>,
%              fillcolor="#DFEDE0", color="#8FB79A", fontcolor="#2E5A3D"];
% 
%   Estimator [label=<<b>Difficulty estimator</b><br/><font point-size="8.5" color="#41719C">Training</font>>,
%              fillcolor="#D3E3F3", color="#7CA6CE", fontcolor="#1F4E79"];
% 
%   Distill   [label=<<b>Distillation</b><br/><font point-size="8.5" color="#41719C">Smaller model</font>>,
%              fillcolor="#D3E3F3", color="#7CA6CE", fontcolor="#1F4E79"];
% 
%   Traffic   -> Agreement [label="Routing signal"];
%   Traffic   -> Corpus    [label="Supervision targets"];
%   Agreement -> Estimator;
%   Corpus    -> Distill;
% 
%   { rank=same; Agreement; Corpus;   }
%   { rank=same; Estimator; Distill;  }
% }
% ```
% label=fig:training_pipeline
% caption=The training pipeline built from one logged dataset.
% rendered_images:end
% render_images:begin
\begin{figure}[H]
  \includegraphics[width=\linewidth]{05_noesis_server.tex.figs/05_noesis_server.5.png}
  \caption{The training pipeline built from one logged dataset.}
  \label{fig:training_pipeline}
\end{figure}
% render_images:end
